\section{Exploratory analyses (post hoc): where the signal lives}\label{sec:exploratory}

Everything in this section is exploratory. The analyses were requested after unblinding, they are outside every
pre-registered family and they support no claim; they suggest where to look, and the forward test of
Section~\ref{sec:limits} re-tests two of them as pre-registered hypotheses. Earlier looks at the test period by
the authors' AI research assistant, at two of the statistics below (hit rates and mean moves by conviction
band, and the statistics of the 10:00 reading), shaped the design of this exploratory study. Every exploratory
variant is logged as a trial: the ledger grew from \res{ntrials-strat} to \res{ntrials-total} trials.

\paragraph{Earnings days carry the signal.} In a post hoc split, every signal's pre-open IC is concentrated on
the events with a scheduled earnings report (\res{earn-n-test} of the \res{test-d0-events} test events), and
there it held in all three periods. In this exploratory split, Laya-Fin's IC with $z_{cc}$ on earnings days was
\res{lf-earn-vala}, \res{lf-earn-valb} and \res{lf-earn-test} on val-A, val-B and the test
($t = \res{lf-earn-test-t}$ on the test), against \res{lf-rest-vala}, \res{lf-rest-valb} and \res{lf-rest-test}
on the other events, and d1-ultra's was \res{d1-earn-vala}, \res{d1-earn-valb} and \res{d1-earn-test}
($t = \res{d1-earn-test-t}$), against \res{d1-rest-vala}, \res{d1-rest-valb} and \res{d1-rest-test}. Post hoc, the
fall from validation to test of Section~\ref{sec:validation} is therefore a fall on the other events. The
exploratory split also shows the commercial models pricing earnings news at the open: d1-ultra's earnings-day IC
on the test is \res{d1-earn-on-test} with $z_{on}$ and \res{d1-earn-id-test} with $z_{id}$.

\paragraph{Selective prediction: the most confident calls carry most of the signal, in the overnight gap.} We
ranked events post hoc by conviction, the size of $|\Nhat|$ against its val-A distribution, and kept only the
top share of them, which traces a coverage curve \citep{geifman2017} (Figure~\ref{fig:selective}). In this
exploratory analysis, d1-ultra's mean overnight move in the direction of its reading,
$\mathrm{sign}(\Nhat)\cdot z_{on}$, is \res{sel-d1-on-10-vala}, \res{sel-d1-on-10-valb} and
\res{sel-d1-on-10-test} typical moves in its top 10\% (val-A, val-B and the test), against
\res{sel-d1-on-all-vala}, \res{sel-d1-on-all-valb} and \res{sel-d1-on-all-test} for all events; for the
close-to-close move the post hoc figures are \res{sel-d1-cc-10-vala}, \res{sel-d1-cc-10-valb} and
\res{sel-d1-cc-10-test} against \res{sel-d1-cc-all-vala}, \res{sel-d1-cc-all-valb} and
\res{sel-d1-cc-all-test}. Post hoc hit rates move less: d1-ultra's top 10\% gets the sign of $z_{cc}$ right
\res{sel-d1-hit-10-test} of the time on the test, against \res{sel-d1-hit-all-test} for all events. After the
open, selectivity gains little in this exploratory analysis: the top 10\%'s mean move in the direction of the
reading is \res{sel-d1-id-10-test} for d1-ultra (95\% interval \res{sel-d1-id-10-test-lo} to
\res{sel-d1-id-10-test-hi}) and \res{sel-lf-id-10-test} for Laya-Fin on the test, and no test interval of this
post hoc mean excludes zero at any coverage.

\begin{figure}[tbp]
\centering
\includegraphics{\figdir/figA16_selective_coverage.pdf}
\caption{\textbf{Exploratory (post hoc): d1-ultra's most confident tenth of pre-open readings carries most of
the overnight signal in all three periods (mean signed overnight move \res{sel-d1-on-10-vala},
\res{sel-d1-on-10-valb} and \res{sel-d1-on-10-test} typical moves, against \res{sel-d1-on-all-vala},
\res{sel-d1-on-all-valb} and \res{sel-d1-on-all-test} for all events), but not the move after the open.}
Coverage curves: events ranked by conviction $|\Nhat|$ (lines) or by the stated confidence (circles), keeping the
most confident share, with thresholds from val-A; 95\% intervals from the stationary bootstrap over sessions.
(a) Mean $\mathrm{sign}(\Nhat)\cdot z_{on}$; (b) mean $\mathrm{sign}(\Nhat)\cdot z_{id}$; (c) the hit rate
against $z_{cc}$. The dashed line marks the frozen confidence gate (the val-A 70th percentile). Requested after
unblinding; the authors' assistant had looked at the conviction bands on the test before this study, which
shaped it. d1-ultra, span-02$^{\ddagger}$ and drex-v2.0 are model-specific.}\label{fig:selective}
\end{figure}

\paragraph{Watching the open adds nothing.} In a post hoc comparison, the reading at 10:00, after the first
half-hour bar, never predicts the rest of the day better than the pre-open reading on val-B or the test (the one
exploratory exception is d1-ultra on val-A, \res{watch-d1-ten-vala} against \res{watch-d1-pre-vala}, neither
significant). Post hoc, the revision from the pre-open to the 10:00 reading has a zero or negative IC with the
move from 10:00 to the close on val-B and the test (Laya-Fin \res{watch-lf-rev-valb} and \res{watch-lf-rev-test},
$t = \res{watch-lf-rev-test-t}$ on the test; d1-ultra \res{watch-d1-rev-valb} and \res{watch-d1-rev-test}): in
this exploratory comparison the models over-update toward the first bar. The comparison involves many
statistics, so a few large $t$-values are expected by chance and the post hoc pattern is weak evidence.

\paragraph{Validation-to-test decay, the strongest instance.} A post hoc conditioning study asked whether the
after-the-open signal $\mathrm{sign}(\Nhat)\cdot z_{id}$ is stronger under some condition: the type and timing
of the news, d1-ultra's own labels of it, the day, the stock, the models' internals, agreement between models
and a cost hurdle, for \res{cond-pairs} (model, condition) pairs. A strict rule of this exploratory study, fixed
before its test run, selected the pairs positive with $t \ge 1$ in both validation quarters and a pooled
$t \ge 2$: \res{cond-selected} pairs passed, with pooled $t$ from \res{cond-pooled-t-min} to
\res{cond-pooled-t-max}. On the test, \res{cond-reversed} of the post hoc selections changed sign and the other
\res{cond-positive} are small (largest $t = \res{cond-pos-t-max}$); in this exploratory family the smallest
Benjamini--Hochberg $q$ is \res{cond-q-selected-min} among the selected pairs and \res{cond-q-min} overall
(Appendix Figure~\ref{fig:a-conditioning}). The exploratory tradable versions had in-sample gross Sharpe ratios
of \res{cb-lf-valb-min} to \res{cb-lf-valb-max} (Laya-Fin) and \res{cb-d1-valb-min} to \res{cb-d1-valb-max}
(d1-ultra) on val-B; on the test these post hoc books' gross Sharpe ratios run from \res{cb-test-min} to
\res{cb-test-max}, none reaches the 90th percentile of its placebo, and every one loses net of 10 basis points.
A post hoc moving-average-cross overlay chosen on validation failed the same way: after crosses of the 20- and
50-day averages without news, the mean abnormal return in the direction of the cross over the next five days was
\res{ma-val-bp} basis points on validation ($t = \res{ma-val-t}$) and \res{ma-test-bp} on the test
($t = \res{ma-test-t}$), so this post hoc reversal did not replicate.
